\documentclass[a4paper, 12pt]{article}

\usepackage{utils}

\renewcommand*{\today}{30 January 2026}

\begin{document}

\hotbox{Analyse 4}{CM 5}{\today}

\begin{definition}
    On dira que la série $\sum\limits_k f_k$ est uniformément de Cauchy sur $I$ si et seulement si
    la suite des sommes partielles $(S_n)$ est uniformément de Cauchy sur $I$.

    C'est-à-dire,
    $$
    \forall \varepsilon > 0, \exists N \in \N, \forall p,q \geq N, \forall x \in I, \norm{S_p - S_q}_{\infty, I} \leq \varepsilon.
    $$
\end{definition}

\begin{remarque}
    Sans perte de généralité, on peut supposer \text{$p \geq q$} et écrire
    $$
    \forall \varepsilon > 0, \exists N \in \N, \forall p \geq q \geq N, \forall x \in I, |S_p(x) - S_q(x)| = \left|\sum_{k=q+1}^p f_k(x)\right| \leq \varepsilon.
    $$
\end{remarque}

\begin{proposition}{}{}
    La série $\sum\limits_k f_k$ converge uniformément sur $I$ si et seulement si elle est uniformément de Cauchy sur $I$.
\end{proposition}

\begin{demonstration}
    \begin{description}
        \item[Sens direct.] Supposons la série $\sum\limits_k f_k$ uniformément convergente vers $S$ \text{sur $I$}.
        
        Soit $\varepsilon > 0$, il existe $N \in \N$ tel que pour tout $n \geq N$,
        $$
        \norm{S_n - S}_{\infty, I} \leq \dfrac{\varepsilon}{2}.
        $$
        Soit $p, q \geq N$, pour tout $x \in I$,
        \begin{align*}
            |S_p(x) - S_q(x)| &= |S_p(x) - S(x) + S(x) - S_q(x)| \\
            &\leq |S_p(x) - S(x)| + |S(x) - S_q(x)| \\
            &\leq \norm{S_p - S}_{\infty, I} + \norm{S_q - S}_{\infty, I} \\
            &\leq \dfrac{\varepsilon}{2} + \dfrac{\varepsilon}{2} = \varepsilon.
        \end{align*}
        Ainsi, la série $\sum\limits_k f_k$ est uniformément de Cauchy sur $I$.

        \item[Sens réciproque.] Supposons que la série $\sum\limits_k f_k$ est uniformément de Cauchy sur $I$.
        
        Soit $x \in I$, la série numérique $\sum\limits_k f_k(x)$ est de Cauchy donc converge simplement vers $S(x)$ (car $\R$ est complet).
        
        Soit $\varepsilon > 0$, il existe $N \in \N$ tel que pour tout $p, q \geq N$, pour tout $x \in I$,
        $$
        |S_p(x) - S_q(x)| \leq \varepsilon.
        $$
        En faisant tendre $q$ vers $+\infty$, on obtient
        $$
        |S_p(x) - S(x)| \leq \varepsilon.
        $$
        Ainsi, la série $\sum\limits_k f_k$ converge uniformément sur $I$.
    \end{description}
\end{demonstration}

\begin{proposition}{}{}
    Si la série $\sum\limits_k f_k$ converge uniformément sur $I$, alors la suite des termes généraux $(f_k)$ converge uniformément vers $\theta$ sur $I$.
\end{proposition}

\begin{demonstration}
    Si la série $\sum\limits_k f_k$ converge uniformément sur $I$, alors la série est uniformément de Cauchy sur $I$.
    
    Il existe $N \in \N$ tel que pour tout $p, q \geq N$, pour tout $x \in I$,
    $$
    \norm{S_p - S_q}_{\infty, I} \underset{p, q \to +\infty}{\to} 0
    $$
    En particulier $q = p + 1$, on a
    $$
    \norm{S_p - S_q}_{\infty, I} = \norm{S_p - S_{p+1}}_{\infty, I} = \norm{f_{p+1}}_{\infty, I} \underset{p \to +\infty}{\to} 0.
    $$
\end{demonstration}

\begin{remarque}
    Le côté intéressant est la contraposée, si \text{$\norm{f_k}_{\infty, I} \not\to 0$} lorsque $k \to +\infty$, alors la série ne converge pas uniformément \text{sur $I$}.
\end{remarque}

\begin{definition}[(Convergence normale)]
    On dira que la série $\sum\limits_k f_k$ converge normalement sur $I$ si et seulement si
    la série numérique $\sum\limits_k \norm{f_k}_{\infty, I}$ converge.
\end{definition}

\begin{methode}
    $\sum\limits_k f_k$ converge normalement sur $I$ si il existe une suite réelle $(\alpha_n)$ telle que
    $\norm{f_k}_{\infty, I} \leq \alpha_k$ et $\sum\limits_k \alpha_k$ converge.
\end{methode}

\begin{theoreme}{}{}
    Si $\sum\limits_k f_k$ converge normalement sur $I$, alors elle converge uniformément sur $I$.
\end{theoreme}

\begin{demonstration}
    Soit $\varepsilon > 0$.

    On sait que $\sum\limits_k f_k$ converge normalement sur $I$, ainsi $\sum\limits_k \norm{f_k}_{\infty, I}$ converge donc la suite de ses sommes partielles est de Cauchy.
    
    Donc il existe $N \in \N$ tel que pour tout $p > q \geq N$,
    $$
    |\sum\limits_{k=k_0}^p \norm{f_k}_{\infty, I} - \sum\limits_{k=k_0}^q \norm{f_k}_{\infty, I}| = \sum\limits_{k=q+1}^p \norm{f_k}_{\infty, I} \leq \varepsilon
    $$

    De plus, on a,

    $$
    \left| \sum_{k=q+1}^p f_k(x) \right| \leq \sum_{k=q+1}^p |f_k(x)| \leq \sum_{k=q+1}^p \norm{f_k}_{\infty, I} \leq \varepsilon
    $$

    car pour tout $x \in I$, $|f_k(x)| \leq \sup\limits_{x \in I}(|f_k(x)|) = \norm{f_k}_{\infty, I}$.

    Ainsi, la série $\sum\limits_k f_k$ est uniformément de Cauchy sur $I$ donc converge uniformément sur $I$.
\end{demonstration}

\begin{exemple}
    Pour tout $n \in \N$, on définit $\funcdef{f_n}{\R}{\R}{x}{\dfrac{(-1)^n}{n + x^2}}$.

    \begin{align*}
        \norm{f_n}_{\infty, \R} &= \sup_{x \in \R}|f_n(x)| \\
        &= \sup_{x \in \R} \left| \dfrac{(-1)^n}{n + x^2} \right| \\
        &= \dfrac{1}{n}
    \end{align*}
    or, la série numérique $\sum_n \dfrac{1}{n}$ diverge.

    Soit $x \in \R, |R_n(x)| = \left| \sum_{k=n+1}^\infty \dfrac{(-1)^n}{n + x^2} \right| \leq \dfrac{1}{n + 1 + x^2} \leq \dfrac{1}{n}$.

    Ainsi, la suite de fonctions des restes $(R_n)$ converge uniformément vers $0$ sur $\R$ donc la série de fonctions $\sum\limits_n f_n$ converge uniformément sur $\R$.
\end{exemple}

\begin{methode}
    Si la série est alternée, on peut utiliser le critère d'Abel
\end{methode}

\begin{exemple}
    Pour tout $n \in \N^*$, on définit $\funcdef{f_n}{\R}{\R}{x}{\dfrac{cos(nx)}{n}}$.

    Il n'y a pas de convergence normale car $\norm{f_n}_{\infty, \R} = \dfrac{1}{n}$ et $\sum\limits_n \dfrac{1}{n}$ diverge.

    Si $x = 2k\pi, k \in \Z, f_n(x) = \dfrac{1}{n}$ donc la série diverge en ces points.

    Si $x \neq 2k\pi, k \in \Z$, on pose $u_n = \dfrac{1}{n}$ et $v_n = \cos(nx)$.

    \begin{align*}
        \sum_{k=1}^n \cos(kx) &= \sum_{k=1}^n Re(e^{ikx}) \\
        &= Re\left( \sum_{k=1}^n (e^{ix})^k \right) \\
        &= Re\left( e^{ix} \dfrac{1 - (e^{ix})^n}{1 - e^{ix}} \right) \\
        &= Re\left( e^{ix} \dfrac{e^{i\dfrac{nx}{2}}\left( e^{-i\dfrac{nx}{2}} - e^{i\dfrac{nx}{2}} \right)}{e^{i\dfrac{x}{2}} \left( e^{-i\dfrac{x}{2}} - e^{i\dfrac{x}{2}} \right)} \right) \\
        &= Re\left( e^{i\dfrac{(n+1)x}{2}} \dfrac{2i \sin(\dfrac{nx}{2})}{2i \sin(\dfrac{x}{2})} \right) \\
        &= \dfrac{\cos(\dfrac{n+1}{2}x) \sin(\dfrac{nx}{2})}{\sin(\dfrac{x}{2})}
    \end{align*}
    ainsi $ \left| V_n \right| \leq \dfrac{1}{\left| \sin(\dfrac{x}{2}) \right|}$.

    Donc par le critère d'Abel, la série $\sum\limits_n f_n$ converge simplement sur tout intervalle $I$ ne contenant pas de $2k\pi, k \in \Z$.

    et on a $|R_n(x)| \leq \dfrac{2}{n \sin(\dfrac{x}{2})}$.

    Par contre $\forall \varepsilon > 0$ et $\varepsilon < k \pi$, on aura convergence uniforme sur $\R_\varepsilon := \bigcup_{k \in \Z}[2k\pi + \varepsilon, 2(k+1)\pi - \varepsilon]$.

    Ainsi, $\forall x \in \R_\varepsilon, |R_n(x)| \leq \dfrac{2}{n \sin(\dfrac{\varepsilon}{2})} \to 0$ lorsque $n \to +\infty$ et ne dépend pas de $x$.

    Donc la série converge uniformément sur $\R_\varepsilon$.
\end{exemple}

\begin{methode}
    On a $1 \pm e^{i\theta} = e^{i\dfrac{\theta}{2}} \left( e^{-i\dfrac{\theta}{2}} \pm e^{i\dfrac{\theta}{2}} \right)$

    donc $1 - e^{i\theta} = 2i \sin(\dfrac{\theta}{2})$ et $1 + e^{i\theta} = 2\cos(\dfrac{\theta}{2})$.

    $\blacktriangle$: Dans le cas où on parle de la phase et le module, si $2\cos(\dfrac{\theta}{2}) < 0$ ou $2i\sin(\dfrac{\theta}{2}) < 0$, il faut ajouter $\pi$ à la phase.
\end{methode}

\begin{theoreme}{Theorème d'inversion des limites}{}
    $\sum_k f_k$ définie sur $I = [a, b]$ et converge uniformément vers $S$ sur $I$.

    On suppose que pour tout $n \in \N$, $f_n(x)$ admet une limite en $b$, notée $\ell_n$.

    Alors $S$ admet une limite en $b$, notée $\ell$ et $(\ell_n)$ converge vers $\ell$ et
\end{theoreme}

\begin{remarque}
    On peut écrire
    $$
    \lim\limits_{x \to b^-} \left( \lim\limits_{n \to +\infty} S_n(x) \right) = \lim\limits_{n \to +\infty} \left( \lim\limits_{x \to b^-} S_n(x) \right)
    $$
\end{remarque}

\begin{theoreme}{Théorème de continuité}{}
    $\sum_k f_k$ définie sur $I$.

    On suppose que $\forall a \in I$, $f_k$ est continue en $a$ et $\sum_k f_k$ converge uniformément vers $S$ sur $I$.

    Alors $S$ est continue sur $I$.
\end{theoreme}

\begin{theoreme}{Théorème de dérivation}{}
    Soit $\sum_k f_k$ définie sur $I$.

    On suppose que pour tout $k \in \N$, $f_k$ est $\mathcal{C}^1$ sur $I$, que $\sum_k f_k$ converge simplement vers $S$ sur $I$ et que $\sum_k f_k'$ converge uniformément sur tout segment de $I$.

    Alors $S$ est $\mathcal{C}^1$ sur $I$ et $S' = \sum_k f_k'$.
\end{theoreme}

\begin{proposition}{}{}
    Soit $\sum_k f_k$ définie sur $I$.

    On suppose que les $f_n$ sont de classe $\mathcal{C}^k$ sur $I$ et que
    \begin{itemize}[label=--]
        \item Pour tout $0 \leq j \leq k - 1$, la série $\sum_n f_n^{(j)}$ converge simplement sur $I$.
        \item La série $\sum_n f_n^{(k)}$ converge uniformément vers $S_k$ sur tout segment de $I$.
    \end{itemize}
    Alors $S = \sum_{l = l_0}^{+\infty} f_l(x)$ est de classe $\mathcal{C}^k$ sur $I$ et pour tout $0 \leq j \leq k$, on a $S^{(j)} = S_j$.
\end{proposition}

\begin{theoreme}{Theorème de permutation somme intégrale}{}
    Soit $\sum_k f_k$ définie sur $I := [a, b]$.

    On suppose que pour tout $n \in \N$, $f_n$ est continue sur $I$ et que $\sum_k f_k$ converge uniformément vers $S$ sur $I$.

    Alors
    $$
    \sum_k \int_a^b f_k(t) \ dt = \int_a^b \sum_k f_k(t) \ dt
    $$
\end{theoreme}

\end{document}